%Deutsch
\newcommand{\hsmaabstractde}{Sprachassistenzsysteme ermöglichen eine komfortable Interaktion mit digitalen Diensten, werfen jedoch aufgrund ihrer häufig cloudbasierten Verarbeitung grundlegende Datenschutzfragen auf. Insbesondere die Erfassung und Verarbeitung von Sprachdaten in privaten Umgebungen macht deutlich, dass alternative Architekturen erforderlich sind, die stärker auf Datenminimierung und lokale Kontrolle ausgerichtet sind. Diese Arbeit untersucht, inwiefern ein lokal betriebener Sprachassistent eine datenschutzfreundliche Alternative zu cloudbasierten Systemen darstellen kann. Dafür wird auf Grundlage der Design Science Research Methodology ein prototypisches System konzipiert, umgesetzt und evaluiert. Der Prototyp kombiniert Wake-Word-Erkennung, Speech-To-Text, lokale \gls{llm}-Inferenz und Text-To-Speech in einer modularen Python-Architektur und verarbeitet die gesamte Interaktion ohne Cloud-Anbindung. Die Evaluation zeigt, dass ein vollständig lokaler Sprachassistent mit heutigen Open-Source-Komponenten grundsätzlich realisierbar ist. Gleichzeitig hängt die praktische Nutzbarkeit stark von der eingesetzten Hardware ab. Während leistungsschwächere Geräte deutliche Einschränkungen bei der Antwortzeit zeigen, ermöglichen aktuelle Systeme eine sehr flüssige Nutzung. Damit zeigt die Arbeit, dass Privacy by Design bei Sprachassistenzsystemen technisch umsetzbar ist, jedoch weiterhin Anforderungen an Ressourcen, Modellwahl und Systemoptimierung stellt.}